% =======================================================================================
% Chapter 1 -- Introduction
% SECTION-H(10): rationale, problem identification, objectives, expected outcomes.
% Chapters 1 and 2 together must not exceed ten pages.
% =======================================================================================
\chapter{Introduction}
\label{ch:introduction}

\section{Rationale of the Research}
\label{sec:rationale}

Speech carries two kinds of information at once. It carries \emph{what} was said, and it carries
\emph{how} it was said --- who the speaker is, how fast they talk, where they place emphasis, how
their pitch rises at the end of a question. Machine translation systems have become very good at
the first and have historically discarded the second. A translated utterance played back in a flat,
unfamiliar synthetic voice is understandable, but it is no longer recognisably the same person
speaking.

Two developments have changed this. Zero-shot voice cloning systems such as
YourTTS \cite{casanova2022yourtts} and XTTS-v2 \cite{casanova2024xtts} synthesise speech in a
target language from a few seconds of reference audio by a speaker who never spoke that language;
and robust multilingual recognition \cite{radford2023whisper} together with openly released
translation models \cite{tiedemann2020opusmt} make a complete speech-to-speech translation (S2ST)
pipeline assemblable on ordinary hardware. The applications are immediate: automatic dubbing of
film and educational video, cross-lingual conferencing, and assistive communication.

The claim usually made for such systems is that they \emph{preserve speaker identity across
languages}. That claim is normally supported with a cosine similarity between speaker embeddings of
the source recording and the synthesised output. The difficulty is that such a number has no natural
scale. A cosine of $0.42$ could indicate excellent preservation or poor preservation; nothing in the
measurement says which, because there is no reference point for how similar the \emph{same person}
sounds to themselves when they switch language. Human speakers do not sound identical across
languages either --- phonetic inventory, rhythm class and habitual pitch range all shift --- so
part of the loss attributed to the synthesiser may simply be the cost of changing language at all.

A second and larger difficulty concerns prosody. Identity and prosody are routinely conflated
under the single heading of ``naturalness'', yet they are separate properties with separate
mechanisms. A voice-cloning synthesiser is conditioned on the reference recording for
\emph{timbre}. It is not, in general, given the source speaker's pitch contour, and it therefore
invents its own intonation from the text it is asked to read. Whether the resulting prosody bears
any relation to what the speaker would actually have done in the target language is an empirical
question, and it can only be answered if a recording of that speaker doing precisely that exists.

This dissertation is built around a corpus for which such recordings do exist. The Dialogs
Re-enacted Across Languages (DRAL) corpus \cite{ward2023dral} asks bilingual participants to hold a
natural conversation in one language and then to re-enact the same exchange, fragment by fragment,
in the other. The result is a set of utterance-aligned pairs in which \emph{one} speaker produces
both sides. That design turns two otherwise unanswerable questions into measurements: how similar
does a speaker sound to themselves across languages, and how much of their prosody carries across.
Both become calibrated ceilings against which a synthesis pipeline can be judged.

\section{Problem Identification}
\label{sec:problem}

Existing bilingual and multilingual speech synthesis systems produce intelligible cross-lingual
speech, but the extent to which they preserve the speaker's individuality and expressive behaviour
remains poorly quantified. Three specific gaps motivate this study.

\textbf{Uncalibrated identity metrics.} Speaker-similarity scores are reported as bare cosines
against the source recording. Without a same-speaker cross-language anchor, it is impossible to say
what fraction of the achievable similarity a system actually attains, or how much of the shortfall
is attributable to the synthesiser rather than to the language change itself.

\textbf{Unmeasured prosody transfer.} Cross-lingual prosody transfer is an active research area
\cite{swiatkowski2023dubbing,song2023styles2st}, but progress is impeded by the scarcity of parallel
expressive data. Work in this area consequently reports perceptual scores or within-language
proxies rather than a direct comparison against what the same speaker did in the target language.
The size of the gap between human and synthetic cross-lingual prosody has therefore not been stated
in units that a subsequent system can be measured against.

\textbf{Confounded reporting.} Correlations computed over a pooled set of utterances from many
speakers conflate two different things: the tendency of a low-pitched person to remain low-pitched
in either language, and the transfer of utterance-level expressive variation. The first is speaker
identity; only the second is prosody transfer. Pooled statistics that do not separate them
systematically overstate how much prosody a system carries.

A fourth, practical problem emerged during the study and is reported alongside the others.
Synthesised utterances are systematically longer than the speech they replace, so a dubbed
recording drifts out of synchronisation with the picture within a few utterances unless the
generated audio is retimed.

\section{Research Questions}
\label{sec:questions}

\begin{enumerate}[label=RQ\arabic*., leftmargin=3.2em]
  \item How much of a speaker's identity, measured against a same-speaker cross-language ceiling,
        survives zero-shot voice cloning in an English--Spanish S2ST cascade?
  \item How much of a speaker's prosody --- pitch level, pitch range, duration and speaking rate
        --- transfers across languages when a human does it, and how much when the system does it?
  \item Is the target-language prosody of an utterance predictable from its source-language
        prosody beyond a trivial copy-plus-offset baseline?
  \item What is the cost, in translation quality, of placing an automatic recogniser in front of
        the translation model, and does in-domain fine-tuning of the translation model recover it?
\end{enumerate}

\section{Aim and Objectives}
\label{sec:objectives}

\textbf{Aim.} To develop and rigorously evaluate a bilingual (English--Spanish) voice transcription
and synthesis system that translates speech between the two languages while preserving the
speaker's vocal identity, and to quantify how much of the speaker's prosody the system carries,
using a same-speaker parallel corpus as the reference standard.

\textbf{Objectives.} The three objectives registered in the approved research proposal, restated
in the terms in which they were finally executed, are:

\begin{enumerate}[label=\textbf{O\arabic*.}, leftmargin=3.0em]
  \item To design and implement a bidirectional voice-to-voice translation framework for English
        and Spanish that preserves the speaker's voice characteristics, using pre-trained ASR, NMT
        and voice-cloning TTS components, and to measure identity preservation against a human
        cross-language reference.
  \item To analyse the prosodic relationship between the two languages for the same speaker, to
        measure how much of that relationship the system reproduces, and to test whether
        target-language prosody can be predicted from source-language prosody.
  \item To integrate the synthesised speech with the original video, applying temporal alignment
        so that the dubbed track keeps the timing of the original, and to deliver this as a working
        application.
\end{enumerate}

A fourth activity, not registered as an objective but required to answer RQ4, was the fine-tuning
of the translation component on the in-domain conversational data and its comparison against the
pre-trained baseline.

\section{Scope and Delimitations}
\label{sec:scope}

The study is deliberately bounded, and the boundaries are stated here so that the results are read
correctly.

\begin{itemize}
  \item \textbf{Two languages, both directions.} English and Spanish only, evaluated in both
        \ento{} and \esto{}.
  \item \textbf{One corpus.} All experimental data come from DRAL. The pre-trained components
        carry knowledge from their own training corpora, which were not modified.
  \item \textbf{Objective evaluation only.} No listening study was conducted. Every result in
        Chapter~\ref{ch:results} is an instrumental measurement. Mean opinion scores, and the
        PESQ and STOI metrics named in the proposal, are not reported; Section~\ref{sec:limitations}
        explains why the latter two are not defined for this task.
  \item \textbf{Analysis, not control.} The prosody model developed in this study
        \emph{predicts} target-language prosody; it does not condition the synthesiser. Closing the
        loop is identified as future work.
  \item \textbf{Duration-matched dubbing, not lip synchronisation.} The application aligns each
        synthesised segment to the duration of the segment it replaces. It does not perform
        phoneme-level lip synchronisation, speaker diarisation, or separation of background music.
\end{itemize}

\section{Expected Outcomes and Contributions}
\label{sec:contributions}

The dissertation delivers:

\begin{enumerate}
  \item A reproducible bidirectional English--Spanish S2ST pipeline with zero-shot voice cloning,
        evaluated on 870 held-out utterances over a speaker-disjoint split, every failure retained.
  \item A \emph{calibrated} identity result --- speaker similarity as a percentage of the human
        same-speaker cross-language ceiling --- and a quantification of the prosody transfer gap
        reported within-speaker, so that identity is not counted twice as prosody.
  \item A negative result with a stated mechanism: target-language prosody is predictable beyond a
        copy-plus-offset baseline only for pitch \emph{range}, and the learned coefficient explains
        why.
  \item A methodological finding: correlating pitch in hertz without a plausibility gate understated
        the correspondence by roughly a factor of seven (Section~\ref{sec:f0correction}).
  \item A working web application that dubs an uploaded video into the other language in the
        original speaker's voice, with each segment retimed to its original slot.
\end{enumerate}

\section{Structure of the Dissertation}
\label{sec:structure}

Chapter~\ref{ch:literature} reviews cascaded and direct S2ST, identity-preserving synthesis,
cross-lingual prosody transfer and the evaluation methodology on which this study depends, and
identifies the gap it addresses. Chapter~\ref{ch:methodology} describes the corpus, the
partitioning, the architecture, every measurement instrument and each experiment.
Chapter~\ref{ch:results} reports and discusses the results in the same order, and
Chapter~\ref{ch:conclusions} evaluates them against the objectives and sets out the next steps.
